\section{The extended statistic and the adjoint method}
\label{ch:ift}

\subsection{Why an implicit derivative is needed}

With externalities or congestion, the welfare effect of a link improvement is
not exhausted by its direct resource saving. The improvement changes realized
costs, routes, modal allocation, bilateral trade, wages, and the location of
labor. These outcomes are determined jointly. Differentiating one equation at
a time while holding the remaining equilibrium objects fixed would omit the
feedback that motivates the exercise.

Collect the transformed equilibrium variables in \(z\) and write the residual
system as
\[
  G(z,\theta)=0.
\]
The vector \(\theta\) contains primitive edge-mode log costs. Around an
interior differentiable baseline,
\[
 J\,dz+G_\theta\,d\theta=0,
 \qquad
 J=G_z.
\]
If \(J\) is nonsingular, the Implicit Function Theorem gives
\[
 \frac{dz}{d\theta}=-J^{-1}G_\theta.
\]
Let \(q^\top=d\log W/dz\) be the welfare row. For a policy \(p\),
\[
 -\frac{d\log W}{d\theta_p}
 =q^\top J^{-1}b_p,
 \qquad b_p=G_{\theta_p}.
\]
The sign convention makes a reduction in cost a positive benefit when the
reported elasticity is positive.

\subsection{The adjoint solve}

There may be hundreds or thousands of candidate links but only one welfare
criterion. Solving \(Jx_p=b_p\) separately for each policy repeats most of the
linear algebra. The adjoint method instead solves
\[
 J^\top a=q
\]
once and evaluates
\[
 E_p=a^\top b_p
\]
for each policy. The result is identical to the direct IFT calculation. After
the adjoint solve, each policy evaluation requires only sparse matrix products.

The package never forms an explicit matrix inverse. It factorizes the relevant
systems and uses linear solves. This matters for numerical stability and for
moderately large sparse networks.

\subsection{Spatial and transport blocks}

The equilibrium has a sparse spatial block and a route-modal-congestion block.
For closure \(S\), write the transport response as
\[
 x=Q_z^S z+C^S a,\qquad
 a=\theta+\Gamma^Sx.
\]
Eliminating the congestion state gives the reduced spatial Jacobian
\[
 J_S
 =J_0+B_\kappa S_{\mathrm{agg}}
 (I-\Gamma^SC^S)^{-1}
 \Gamma^SQ_z^S.
\]
Here:
\begin{itemize}
  \item \(J_0\) is the no-congestion spatial Jacobian;
  \item \(B_\kappa\) loads aggregate edge costs into spatial residuals;
  \item \(S_{\mathrm{agg}}\) aggregates edge-mode costs to physical directed
  edges;
  \item \(Q_z^S\) maps the spatial state into edge-mode traffic;
  \item \(C^S\) maps congestion states into traffic;
  \item \(\Gamma^S\) maps traffic into realized costs.
\end{itemize}

This Schur complement is not a numerical shortcut detached from the theory. It
is the result of eliminating the differentiated transport quantities from the
uneliminated residual system. Tests compare the two systems directly.

\subsection{Realized and primitive shocks}

Two derivatives must be distinguished. A realized-friction shock changes the
cost users face after congestion:
\[
 \vartheta_{klm}=\log\kappa_{kl,m}.
\]
A primitive policy shock changes the infrastructure component:
\[
 \theta_{klm}=\log\bar\kappa_{kl,m}.
\]
With congestion, these are not the same experiment. Traffic induced by a
primitive improvement can attenuate the fall in realized cost.

Let
\[
 \zeta_{klm}
 =\left.
 \frac{\partial\log\kappa_{kl}}
      {\partial\theta_{klm}}
 \right|_{dz=0},
 \qquad
 \chi_{klm}=\frac{\zeta_{klm}}{s_{kl,m}}.
\]
The derivative is taken through the modal and congestion system while holding
the spatial state fixed. The spatial equilibrium then responds through
\(J^{-1}\). The implementation constructs the full primitive forcing rather
than assigning \(\chi\) separately.

When the realized welfare derivative is nonzero, the output reports the
welfare-effective ratio
\[
 \chi_e^{eff}=\frac{\Eprimitive_{e,F}}{\Ereal_{e,F}}.
\]
If the denominator is numerically zero, the ratio is missing and the
pass-through contribution is still reported in levels. A zero derivative is a
valid result, not an exception.

\subsection{The economics-facing proposition}

The algebraic IFT expression can be written as traffic times equilibrium
multipliers. Define
\[
 \rho=\frac{1+\alpha+\beta}
 {1+\beta(\sigma-1)+\alpha\sigma}.
\]
Let \(\mathcal M_i^{in}\) and \(\mathcal M_i^{out}\) be the adjoint
market-access multipliers normalized by trade exposure \(\mathcal T_i\). Then
the primitive edge-mode welfare elasticity is
\[
 -\frac{d\log W}{d\theta_{klm}}
 =\chi_{klm}\rho\Xi_{kl,m}
 \left(\mathcal M_k^{in}+\mathcal M_l^{out}\right).
\]

The terms have distinct roles:
\begin{description}
  \item[Traffic \(\Xi_{kl,m}\).] The scale of economic activity directly
  exposed to the improvement.
  \item[Externality scale \(\rho\).] A common transformation of welfare and
  spatial responses for a given parameter vector.
  \item[Pass-through \(\chi_{klm}\).] The part of the primitive change that
  reaches the aggregate edge cost after modal and congestion responses.
  \item[Market-access multipliers.] The link-specific incidence of the shock
  through inward and outward equilibrium propagation.
\end{description}

At the paper calibration, \(e=-0.5\) and \(\rho=-1.6\). The sign of \(\rho\)
alone does not determine the sign of the welfare effect because the multiplier
sum is also signed. The reported gains remain positive at the accepted
baseline.

\subsection{Observable and constructed objects}

Traffic, labor shares, and mode shares can be observed or constructed directly.
Route shares and trade exposure depend on the recursive network accounting.
Using outgoing or incoming traffic,
\[
 \mathcal T_i=\frac{O_i}{1-s_i^x}
             =\frac{I_i}{1-s_i^y}.
\]
This is an identity in the model. Empirical implementation requires a
model-consistent flow network. The package reports stock disagreement rather
than hiding it through an automatic balancing step.

\implementationbox{\code{build\_transport\_closure} constructs the Schur
complement. \code{primitive\_forcing} differentiates the primitive policy.
\code{operator\_gain} performs the adjoint projection. The sparse
\code{edge\_local\_welfare\_effects} path computes the paper's headline
edge-local application without constructing the dense fixed-route incidence
needed for the full decomposition.}

\subsection{Regularity and failure}

The derivative is local and conditional on:
\begin{itemize}
  \item an interior baseline with positive active edge-mode flows;
  \item a contractive route kernel;
  \item positive exposure stocks;
  \item \(e\neq0\) and \(1+\alpha+\beta\neq0\);
  \item nonsingular modal-congestion and spatial systems;
  \item a condition number below the configured gate.
\end{itemize}
The software fails explicitly when these conditions are violated. It never
uses a ridge term to turn an unidentified derivative into a number.
